%======================================================================
\chapter{Notaci\'on}
\label{ap:notacion}

\begin{longtable}{@{}lp{10.2cm}@{}}
\toprule
S\'imbolo & Significado \\
\midrule
\endfirsthead
\toprule
S\'imbolo & Significado \\
\midrule
\endhead
$\hat\varepsilon=\varepsilon_1+i\varepsilon_2$ & Funci\'on diel\'ectrica compleja \\
$\hat N=n+ik$ & \'Indice de refracci\'on complejo \\
$R$ & Reflectividad a incidencia normal (semi-infinito) \\
$\alpha$ & Coeficiente de absorci\'on \\
$T$ & Transmitancia de una l\'amina \\
$\varepsilon_\infty$ & Fondo de alta frecuencia (lo no medido) \\
$\Lambda_k$, $A_k$ & Tri\'angulo unitario y su amplitud \\
$K_k$ & Transformada Kramers--Kronig de $\Lambda_k$ \\
$\Phi$ & Objetivo regularizado ($\chi^{2}_{\mathrm{datos}}+P_s+P_-$) \\
$\chi^{2}_{\mathrm{datos}}$ & Suma de residuos cuadr\'aticos de la magnitud medida \\
RefFIT & Programa de Kuzmenko para ajuste de espectros \'opticos \\
\bottomrule
\end{longtable}

\chapter{Lista de pruebas A--G}
\label{ap:pruebas}

\begin{enumerate}[leftmargin=*]
\item[A.] Cierre algebraico $\varepsilon_1=n^{2}-k^{2}$, $\varepsilon_2=2nk$.
\item[B.] Proyecci\'on de la medida: $\chi^{2}_{\mathrm{datos}}$ y residuos.
\item[C.] Kramers--Kronig num\'erica independiente.
\item[D.] Sensibilidad a $\varepsilon_\infty$.
\item[E.] Barrido de regularizaci\'on (suavidad).
\item[F.] Recuperaci\'on sint\'etica (tabla~\ref{tab:pruebaF}).
\item[G.] Comparaci\'on cruzada con RefFIT (tabla~\ref{tab:reffit-plantilla}).
\end{enumerate}

\chapter{Plantilla de reporte reproducible}
\label{ap:reporte}

Todo resultado que se presente en un seminario o en un art\'iculo derivado debe incluir, como m\'inimo:

\begin{enumerate}[leftmargin=*]
\item archivo de datos (o hash), unidad de $\omega$ y magnitud medida;
\item $\varepsilon_\infty$ (valor fijo o regla auto) y espesor si aplica;
\item suavidad, negatividad, n\'umero de anclas e iteraciones;
\item $\chi^{2}_{\mathrm{datos}}$, estado de convergencia y aviso de estancamiento si lo hubo;
\item curvas $y$ frente a $Q$, $\varepsilon_1$, $\varepsilon_2$;
\item si hubo Prueba~G: controles id\'enticos en RefFIT y tabla~\ref{tab:reffit-plantilla} completa.
\end{enumerate}

\chapter{Itinerario de seminarios (ULA)}
\label{ap:seminarios}

\begin{description}[style=nextline,leftmargin=1.2cm]
\item[Seminario~I.]
Cap\'itulos~\ref{ch:problema}--\ref{ch:kk}. Mensaje: causalidad y $\hat\varepsilon$.
\item[Seminario~II.]
Cap\'itulos~\ref{ch:variacional}--\ref{ch:app}. Mensaje: tri\'angulos, $\Phi$, cada medida, \textbf{arquitectura Rust} (\texttt{measure}, \texttt{fit}, \texttt{http}, \texttt{run\_fit}).
\item[Seminario~III (final).]
Cap\'itulos~\ref{ch:verdad}--\ref{ch:conclusiones} y Secci\'on~\ref{sem:3}. Mensaje: validaci\'on, RefFIT, l\'imites.
\end{description}

\chapter{Mapa r\'apido del c\'odigo Rust (Seminario~II)}
\label{ap:rust}

\begin{longtable}{@{}lp{9.6cm}@{}}
\toprule
Archivo / s\'imbolo & Rol \\
\midrule
\endfirsthead
\toprule
Archivo / s\'imbolo & Rol \\
\midrule
\endhead
\texttt{measure::quantity} & Proyecta $(\varepsilon_1,\varepsilon_2)\to Q$ (la magnitud medida) \\
\texttt{measure::guess\_eps\_inf} & Auto $\varepsilon_\infty$ desde cola de $R$ o desde $n$ \\
\texttt{fit::kk\_kernel} & N\'ucleo anal\'itico $K_k$ del tri\'angulo \\
\texttt{fit::mesh\_sample\_indices} & Elige anclas (picos/valles + fondo) \\
\texttt{fit::Model} & $\varepsilon_1$, $\varepsilon_2$, $\chi^{2}$ y derivadas en un punto de par\'ametros \\
\texttt{fit::levenberg} & Paso amortiguado, cotas, parada temprana \\
\texttt{fit::run\_fit} & Pipeline completo $\to$ \texttt{FitReport} \\
\texttt{fit::fit\_text\_with\_progress} & Entrada texto + ganchos de progreso \\
\texttt{http::fit\_worker} & Job en servidor, cancelaci\'on, curva en vivo \\
\texttt{store::*} & Disco: jobs e informes \\
\bottomrule
\end{longtable}
